\documentclass{csexam}

\usepackage{listings}
\usepackage{xcolor}
\usepackage{tikz}

% Python code styling
\lstset{
    language=Python,
    basicstyle=\ttfamily\small,
    keywordstyle=\color{blue}\bfseries,
    stringstyle=\color{red},
    commentstyle=\color{green!50!black}\itshape,
    showstringspaces=false,
    breaklines=true,
    frame=single,
    numbers=left,
    numberstyle=\tiny\color{gray}
}

% Course information
\university{St. Francis Xavier University}
\department{Department of Computer Science}
\course{CSCI 128: Coding for Problem Solving}
\courseshort{CSCI 128}
\exam{Midterm 1}
\examdate{Winter 2026}

\begin{document}

\maketitle

\begin{rules}
\textbf{Instructions:}
\begin{itemize}
    \item This exam is \textbf{50 minutes} long.
    \item Answer all questions in the space provided.
    \item This exam is worth \textbf{100 points}.
    \item No calculators, books, or notes are permitted.
    \item Write your answers clearly. Illegible answers will receive no credit.
    \item For coding questions, write syntactically correct Python code.
    \item Good luck!
\end{itemize}
\end{rules}

\vspace{5mm}

\showgrade

\newpage

\begin{questions}

% ======================================
% Section 1: Values, Types, and Variables
% ======================================
\question[6] For each expression below, determine the data type of the result. If there is an error, write "ERROR".

\begin{parts}
    \part \texttt{5 + 3.0} \hspace{0.5cm} Type: \underline{\hspace{3cm}}
    \part \texttt{str(42) + "7"} \hspace{0.5cm} Type: \underline{\hspace{3cm}}
    \part \texttt{int("3.14")} \hspace{0.5cm} Type/Result: \underline{\hspace{3cm}}
    \part \texttt{10 / 2} \hspace{0.5cm} Type: \underline{\hspace{3cm}}
    \part \texttt{10 // 3} \hspace{0.5cm} Type: \underline{\hspace{3cm}}
    \part \texttt{"Python" * 2} \hspace{0.5cm} Type: \underline{\hspace{3cm}}
\end{parts}

\question[8] What is the output of the following code? Show exactly what gets printed.

\begin{lstlisting}
x = "5"
y = 3
x = int(x) + y
y = str(y) * 2
print(x, y)
\end{lstlisting}
\vspace{2.5cm}

\question[8] Trace through the following code carefully and show what gets printed:
\begin{lstlisting}
a = 10
b = 20
a = b - a
b = a + b
a = b - a
print(a)
print(b)
\end{lstlisting}
\vspace{3cm}

\newpage

\question[14] Write a complete Python program that:
\begin{itemize}
    \item Asks the user for their age
    \item Asks the user for the number of years to add
    \item Calculates their future age
    \item Calculates how many months old they will be (future age × 12)
    \item Prints both results with appropriate labels
\end{itemize}

\textbf{Example run:}
\begin{verbatim}
Enter your age: 20
Enter number of years: 5
In 5 years you will be: 25
That is 300 months old
\end{verbatim}

\vspace{8cm}

\newpage

\question[8] What is wrong with each of the following code snippets? Identify the error and explain how to fix it.

\begin{parts}
    \part
\begin{lstlisting}
age = input("Enter your age: ")
next_year = age + 1
print(next_year)
\end{lstlisting}
    \textbf{Error:} \underline{\hspace{10cm}}

    \textbf{Fix:} \underline{\hspace{10cm}}
    \vspace{1cm}

    \part
\begin{lstlisting}
name = input("Enter your name: ")
print("Hello " + name + ". You are " + 25 + " years old")
\end{lstlisting}
    \textbf{Error:} \underline{\hspace{10cm}}

    \textbf{Fix:} \underline{\hspace{10cm}}
    \vspace{1cm}
\end{parts}

\question[6] Explain the difference between the operators \texttt{/} and \texttt{//} in Python. Give an example showing different results for each.
\vspace{4cm}

\newpage

% ======================================
% Section 2: Pictures and Pixels
% ======================================
\question[6] Answer the following questions:

\begin{parts}
    \part What color would result from the RGB values (128, 128, 128)? What about (200, 200, 200)?
    \vspace{1.5cm}

    \part What RGB values would you use to create yellow?
    \vspace{1.5cm}

    \part In the RGB color model, how would you create orange? (Hint: orange is red + some green)
    \vspace{1.5cm}
\end{parts}

\question[5] For an image with dimensions 100×100:

\begin{parts}
    \part What are the coordinates of the top-right corner pixel?
    \vspace{1.5cm}

    \part What are the coordinates of the bottom-left corner pixel?
    \vspace{1.5cm}

    \part If you wanted to access the pixel in the exact center of the image, what coordinates would you use?
    \vspace{1.5cm}
\end{parts}

\newpage

\question[15] Write a complete Python program using Pillow that:
\begin{itemize}
    \item Loads an image file named \texttt{"photo.jpg"}
    \item Gets the width and height of the image and prints them
    \item Gets the pixel color at position (10, 10)
    \item Extracts the individual R, G, and B values from that pixel
    \item Calculates the sum of the three color values (R + G + B)
    \item Prints the results in this format:
    \begin{itemize}
        \item "Image size: [width] x [height]"
        \item "Pixel at (10, 10) - Red: [R], Green: [G], Blue: [B]"
        \item "Total color value: [sum]"
    \end{itemize}
\end{itemize}

\textit{Hint: The pixel color is returned as a tuple (R, G, B). You can unpack it like: r, g, b = pixel\_color}

\vspace{10cm}

\newpage

\question[12] Complete the following Python program by filling in the blanks. This program creates a vertical red stripe down the middle of a 200×200 image:

\begin{lstlisting}
from PIL import Image

# Create a new white image (200x200)
img = Image.______("RGB", (______, ______), (255, 255, 255))

# Draw a red vertical stripe in the middle (x from 95 to 105)
for x in range(______, ______):
    for y in range(______):
        img.______((x, y), (______, ______, ______))

# Save the image
img.______("stripe.jpg")
\end{lstlisting}

\vspace{2cm}

\question[12] Write a complete Python program using Pillow that:
\begin{itemize}
    \item Creates a new 100×100 RGB image with all pixels initially black
    \item Uses nested loops to set all pixels in the left half (x from 0 to 49) to red (255, 0, 0)
    \item Uses nested loops to set all pixels in the right half (x from 50 to 99) to green (0, 255, 0)
    \item Saves the result as \texttt{"two\_colors.jpg"}
\end{itemize}

\textit{Remember: You need nested loops (one for x, one for y) to access all pixels.}

\vspace{10cm}

\end{questions}

\end{document}
